\section{Introduction}

Spaceborne synthetic aperture radar (SAR) provides wide-area maritime
observation independent of daylight and through most weather
conditions~\cite{torres2012sentinel1}. Those properties make Sentinel-1
imagery useful for detecting vessels without a contemporaneous matched
Automatic Identification System report---the so-called \emph{dark vessels}.
Such detections support maritime safety, fisheries monitoring, and
investigation of activity potentially associated with sanctions evasion. The
task is difficult: vessels occupy few pixels in a large scene, sea clutter
and coastal structures produce strong responses, and human-verified labels
are expensive. The xView3-SAR benchmark makes the problem concrete with
ground-range-detected imagery, georeferenced point annotations, and an
evaluation protocol based on geographic matching~\cite{paolo2022xview3}.

The scarcity of verified labels motivates transfer learning, but it is not
clear which source helps most. Generic ImageNet pretraining supplies broad
visual features; optical remote-sensing pretraining adds overhead-scene
statistics; SAR pretraining additionally shares the sensing modality. These
alternatives also interact with architecture: vision transformers and
convolutional networks encode different inductive biases, and released
remote-sensing checkpoints are not available under one method-matched
recipe. A comparison that changes architecture, optimization, detector, and
pretraining source simultaneously cannot isolate the source-domain question.

We study label efficiency through two architecture-matched tracks
(Fig.~\ref{fig:study-design}, Appendix~\ref{app:figures}). Within each
track, four arms differ only in initialization: random, ImageNet, optical
remote sensing, or SAR. Every arm fine-tunes end to end through the same
point detector, input representation, optimizer, schedule, nested scene
subsets, and fixed seed, producing 32 controlled cells. The study asks
three questions:
\begin{enumerate}
\item How does random, ImageNet, optical, or SAR initialization change
  detection F1 as the number of labeled scenes grows?
\item Does SAR pretraining beat optical pretraining within each
  architecture, especially under label scarcity, and is the direction of
  that contrast consistent across tracks?
\item After model selection and thresholds freeze, how do the arms
  generalize to held-out scenes and human-verified dark vessels?
\end{enumerate}
This report answers the first two questions with the completed 32-cell H100
campaign, every value re-derived at build time from hash-bound
evidence.\ifHevTestComplete{} The once-scored 16-scene test matrix addresses
the held-out part of the third.\else{} The 16-scene test matrix and the
sealed 50-scene human-verified evaluation address the third and are
described in Sect.~\ref{sec:results}.\fi{} Our contribution is experimental
control rather than a new method: an eight-arm two-track design in which
initialization is the only within-track variable, one corrected evaluation
contract with checkpoint-bound operating thresholds, and a fail-closed
evidence path from training bytes to every published number.
